\renewcommand{\tema}{Preguntas Post-Clase - Cinemática}

\twocolumn
\section{PREGUNTAS POST-CLASE}

\begin{preguntabox}
\subsection{Pregunta 1}
Un automóvil viaja a 90 km/h durante 20 minutos. ¿Qué distancia recorre?

\begin{enumerate}
    \item[A)] 18 km
    \item[B)] 30 km
    \item[C)] 45 km
    \item[D)] 90 km
\end{enumerate}
\end{preguntabox}

\begin{preguntabox}
\subsection{Pregunta 2}
Un móvil recorre 200 m hacia el este y luego 150 m hacia el oeste. Si el tiempo total fue de 50 s, ¿cuál es la magnitud de su velocidad promedio y su rapidez promedio?

\begin{enumerate}
    \item[A)] Velocidad: 1 m/s, Rapidez: 7 m/s
    \item[B)] Velocidad: 7 m/s, Rapidez: 1 m/s
    \item[C)] Velocidad: 1 m/s, Rapidez: 1 m/s
    \item[D)] Velocidad: 7 m/s, Rapidez: 7 m/s
\end{enumerate}
\end{preguntabox}

\begin{preguntabox}
\subsection{Pregunta 3}
Un tren acelera uniformemente desde 10 m/s hasta 30 m/s en 8 segundos. ¿Cuál es su aceleración?

\begin{enumerate}
    \item[A)] 1.25 m/s²
    \item[B)] 2.5 m/s²
    \item[C)] 3.75 m/s²
    \item[D)] 5 m/s²
\end{enumerate}
\end{preguntabox}

\begin{preguntabox}
\subsection{Pregunta 4}
Un objeto parte del reposo con aceleración de 4 m/s². ¿Qué distancia recorre en los primeros 5 segundos?

\begin{enumerate}
    \item[A)] 10 m
    \item[B)] 20 m
    \item[C)] 50 m
    \item[D)] 100 m
\end{enumerate}
\end{preguntabox}

\newpage
\begin{preguntabox}
\subsection{Pregunta 5}
La siguiente gráfica velocidad vs tiempo muestra el movimiento de un objeto. ¿Qué distancia recorre?

\begin{center}
\begin{tikzpicture}
\begin{axis}[
    width=7cm,
    height=5cm,
    axis lines=middle,
    xlabel={$t$ (s)},
    ylabel={$v$ (m/s)},
    xmin=0, xmax=7,
    ymin=0, ymax=10,
    xtick={0,2,4,6},
    ytick={0,2,4,6,8},
    grid=major,
    thick,
    every axis x label/.style={at={(current axis.right of origin)},anchor=west},
    every axis y label/.style={at={(current axis.above origin)},anchor=south}
]
% Rectángulo: velocidad constante de 8 m/s durante 6 s
\addplot[color=azuloscuro, very thick] coordinates {(0,8) (6,8)};
\addplot[color=azuloscuro, very thick] coordinates {(0,0) (0,8)};
\addplot[color=azuloscuro, very thick] coordinates {(6,0) (6,8)};
\addplot[color=azuloscuro, very thick] coordinates {(0,0) (6,0)};
\end{axis}
\end{tikzpicture}
\end{center}

\begin{enumerate}
    \item[A)] 14 m
    \item[B)] 24 m
    \item[C)] 48 m
    \item[D)] 64 m
\end{enumerate}
\end{preguntabox}

\begin{preguntabox}
\subsection{Pregunta 6}
Un objeto se lanza verticalmente hacia arriba con velocidad de 50 m/s. ¿Qué altura máxima alcanza? Considera $g = 10$ m/s²

\begin{enumerate}
    \item[A)] 25 m
    \item[B)] 50 m
    \item[C)] 125 m
    \item[D)] 250 m
\end{enumerate}
\end{preguntabox}

\begin{preguntabox}
\subsection{Pregunta 7}
Un móvil frena uniformemente desde 20 m/s hasta detenerse, recorriendo 40 m. ¿Cuál fue su aceleración?

\begin{enumerate}
    \item[A)] -10 m/s²
    \item[B)] -5 m/s²
    \item[C)] -2.5 m/s²
    \item[D)] -0.5 m/s²
\end{enumerate}
\end{preguntabox}

\newpage
\begin{preguntabox}
\subsection{Pregunta 8}
Un auto duplica su velocidad en 4 segundos, recorriendo 60 m en ese tiempo. Si partió con velocidad $v_0$, ¿cuál es $v_0$?

\begin{enumerate}
    \item[A)] 5 m/s
    \item[B)] 10 m/s
    \item[C)] 15 m/s
    \item[D)] 20 m/s
\end{enumerate}
\end{preguntabox}

\begin{preguntabox}
\subsection{Pregunta 9}
Se deja caer una piedra desde un edificio. Si tarda 4 segundos en llegar al suelo, ¿cuál es la altura del edificio? Considera $g = 10$ m/s²

\begin{enumerate}
    \item[A)] 40 m
    \item[B)] 60 m
    \item[C)] 80 m
    \item[D)] 160 m
\end{enumerate}
\end{preguntabox}

\begin{preguntabox}
\subsection{Pregunta 10}
En una gráfica posición vs tiempo, un objeto está en $x = 40$ m en $t = 2$ s y en $x = 10$ m en $t = 5$ s. ¿Cuál es su velocidad?

\begin{enumerate}
    \item[A)] -15 m/s
    \item[B)] -10 m/s
    \item[C)] -5 m/s
    \item[D)] 10 m/s
\end{enumerate}
\end{preguntabox}


\begin{preguntabox}
\subsection{Pregunta 11}
Dos móviles parten del mismo punto. El primero a 15 m/s constante, el segundo 4 s después a 25 m/s. ¿A qué distancia del origen se encuentran?

\begin{enumerate}
    \item[A)] 60 m
    \item[B)] 100 m
    \item[C)] 150 m
    \item[D)] 200 m
\end{enumerate}
\end{preguntabox}

\begin{preguntabox}
\subsection{Pregunta 12}
Un objeto con MUA parte del reposo. Si en los primeros 2 s recorre 8 m, ¿cuánto recorre en los primeros 4 s?

\begin{enumerate}
    \item[A)] 16 m
    \item[B)] 24 m
    \item[C)] 32 m
    \item[D)] 64 m
\end{enumerate}
\end{preguntabox}

\begin{preguntabox}
\subsection{Pregunta 13}
Se lanza un proyectil horizontalmente desde 45 m de altura con velocidad de 30 m/s. ¿Qué distancia horizontal recorre? Considera $g = 10$ m/s²

\begin{enumerate}
    \item[A)] 30 m
    \item[B)] 60 m
    \item[C)] 90 m
    \item[D)] 135 m
\end{enumerate}
\end{preguntabox}

\begin{preguntabox}
\subsection{Pregunta 14}
La siguiente gráfica velocidad vs tiempo muestra el movimiento de un móvil que parte del reposo. ¿Qué aceleración tiene?

\begin{center}
\begin{tikzpicture}
\begin{axis}[
    width=7cm,
    height=5cm,
    axis lines=middle,
    xlabel={$t$ (s)},
    ylabel={$v$ (m/s)},
    xmin=0, xmax=9,
    ymin=0, ymax=14,
    xtick={0,2,4,6,8},
    ytick={0,3,6,9,12},
    grid=major,
    thick,
    every axis x label/.style={at={(current axis.right of origin)},anchor=west},
    every axis y label/.style={at={(current axis.above origin)},anchor=south}
]
% Triángulo: de (0,0) a (8,12)
\addplot[color=azuloscuro, very thick] coordinates {(0,0) (8,12)};
\addplot[color=azuloscuro, very thick] coordinates {(0,0) (8,0)};
\addplot[color=azuloscuro, very thick] coordinates {(8,0) (8,12)};
\end{axis}
\end{tikzpicture}
\end{center}

\begin{enumerate}
    \item[A)] 0.67 m/s²
    \item[B)] 1.5 m/s²
    \item[C)] 3 m/s²
    \item[D)] 6 m/s²
\end{enumerate}
\end{preguntabox}

\begin{preguntabox}
\subsection{Pregunta 15}
Un objeto se lanza verticalmente hacia arriba y regresa al punto de lanzamiento en 6 s. ¿Con qué velocidad se lanzó? Considera $g = 10$ m/s²

\begin{enumerate}
    \item[A)] 15 m/s
    \item[B)] 20 m/s
    \item[C)] 30 m/s
    \item[D)] 60 m/s
\end{enumerate}
\end{preguntabox}

\newpage
\onecolumn
\section{RESPUESTAS Y RETROALIMENTACIÓN}

\begin{respuestabox}
\subsection{Pregunta 1}
\textcolor{verdecorrecto}{\textbf{Respuesta correcta: B) 30 km}}

\textbf{Justificación:}\\
\textbf{Fórmula utilizada:} Distancia en movimiento rectilíneo uniforme (Lección: Movimiento Rectilíneo Uniforme)
$$d = vt$$

\textbf{Despeje:} Ya está despejada para $d$

\textbf{Conversión de unidades:}
$$t = 20 \text{ min} \times \frac{1 \text{ h}}{60 \text{ min}} = \frac{1}{3} \text{ h}$$

\textbf{Sustitución:}
$$d = 90 \times \frac{1}{3} = 30 \text{ km}$$

\textbf{Respuestas incorrectas:}
\begin{itemize}
    \item \textbf{A) 18 km:} Error en la conversión de tiempo o en la multiplicación.
    \item \textbf{C) 45 km:} Convierte 20 minutos como 0.5 horas en lugar de 0.33 horas.
    \item \textbf{D) 90 km:} Usa el tiempo como 1 hora completa en lugar de 1/3 de hora.
\end{itemize}
\end{respuestabox}

\begin{respuestabox}
\subsection{Pregunta 2}
\textcolor{verdecorrecto}{\textbf{Respuesta correcta: A) Velocidad: 1 m/s, Rapidez: 7 m/s}}

\textbf{Justificación:}

\textbf{Para velocidad promedio:}\\
\textbf{Fórmula:} Velocidad promedio (Lección: Características de los fenómenos mecánicos)
$$v_{prom} = \frac{\Delta x}{t}$$

\textbf{Cálculo del desplazamiento:}
$$\Delta x = 200 - 150 = 50 \text{ m (hacia el este)}$$

\textbf{Sustitución:}
$$v_{prom} = \frac{50}{50} = 1 \text{ m/s}$$

\textbf{Para rapidez promedio:}\\
\textbf{Fórmula:} Rapidez promedio (Lección: Características de los fenómenos mecánicos)
$$r_{prom} = \frac{d_{total}}{t}$$

\textbf{Cálculo de distancia total:}
$$d_{total} = 200 + 150 = 350 \text{ m}$$

\textbf{Sustitución:}
$$r_{prom} = \frac{350}{50} = 7 \text{ m/s}$$

\textbf{Respuestas incorrectas:}
\begin{itemize}
    \item \textbf{B) Invertidas:} Confunde velocidad con rapidez.
    \item \textbf{C) Ambas iguales a 1:} Usa desplazamiento para ambas.
    \item \textbf{D) Ambas iguales a 7:} Usa distancia total para ambas.
\end{itemize}
\end{respuestabox}

\newpage

\begin{respuestabox}
\subsection{Pregunta 3}
\textcolor{verdecorrecto}{\textbf{Respuesta correcta: B) 2.5 m/s²}}

\textbf{Justificación:}\\
\textbf{Fórmula utilizada:} Definición de aceleración (Lección: Movimiento uniformemente acelerado)
$$a = \frac{v_f - v_i}{t}$$

\textbf{Despeje:} Ya está despejada para $a$

\textbf{Sustitución:}
$$a = \frac{30 - 10}{8} = \frac{20}{8} = 2.5 \text{ m/s}^2$$

\textbf{Respuestas incorrectas:}
\begin{itemize}
    \item \textbf{A) 1.25 m/s²:} Divide entre 16 en lugar de 8.
    \item \textbf{C) 3.75 m/s²:} Suma las velocidades antes de dividir: $\frac{40}{8}$.
    \item \textbf{D) 5 m/s²:} Error en el denominador.
\end{itemize}
\end{respuestabox}

\begin{respuestabox}
\subsection{Pregunta 4}
\textcolor{verdecorrecto}{\textbf{Respuesta correcta: C) 50 m}}

\textbf{Justificación:}\\
\textbf{Fórmula utilizada:} Distancia en MUA (Lección: Movimiento uniformemente acelerado)
$$d = v_i t + \frac{1}{2}at^2$$

\textbf{Despeje:} Ya está despejada para $d$

\textbf{Datos:} $v_i = 0$ (parte del reposo), $a = 4$ m/s², $t = 5$ s

\textbf{Sustitución:}
$$d = 0(5) + \frac{1}{2}(4)(5)^2 = 0 + 2(25) = 50 \text{ m}$$

\textbf{Respuestas incorrectas:}
\begin{itemize}
    \item \textbf{A) 10 m:} Olvida elevar al cuadrado el tiempo.
    \item \textbf{B) 20 m:} Usa $at$ en lugar de $\frac{1}{2}at^2$.
    \item \textbf{D) 100 m:} Olvida el factor $\frac{1}{2}$.
\end{itemize}
\end{respuestabox}

\newpage

\begin{respuestabox}
\subsection{Pregunta 5}
\textcolor{verdecorrecto}{\textbf{Respuesta correcta: C) 48 m}}

\textbf{Justificación:}\\
\textbf{Concepto utilizado:} Área bajo la curva en gráfica v-t (Lección: Movimiento rectilíneo uniforme - interpretación gráfica)

\textbf{Fórmula del área de un rectángulo:}
$$d = \text{base} \times \text{altura}$$

\textbf{Datos de la gráfica:} base = 6 s, altura = 8 m/s

\textbf{Sustitución:}
$$d = 6 \times 8 = 48 \text{ m}$$

\textbf{Respuestas incorrectas:}
\begin{itemize}
    \item \textbf{A) 14 m:} Suma base y altura en lugar de multiplicar.
    \item \textbf{B) 24 m:} Calcula solo la mitad del área.
    \item \textbf{D) 64 m:} Eleva al cuadrado la altura.
\end{itemize}
\end{respuestabox}


\begin{respuestabox}
\subsection{Pregunta 6}
\textcolor{verdecorrecto}{\textbf{Respuesta correcta: C) 125 m}}

\textbf{Justificación:}\\
\textbf{Fórmula utilizada:} Ecuación de velocidad-distancia en MUA (Lección: Movimiento uniformemente acelerado)
$$v_f^2 = v_i^2 + 2ad$$

\textbf{Condición:} En la altura máxima, $v_f = 0$. La aceleración es $a = -g = -10$ m/s²

\textbf{Despeje para $d$ (altura $h$):}
$$v_f^2 = v_i^2 - 2gh$$
$$0 = v_i^2 - 2gh$$
$$2gh = v_i^2$$
$$h = \frac{v_i^2}{2g}$$

\textbf{Sustitución:}
$$h = \frac{(50)^2}{2(10)} = \frac{2500}{20} = 125 \text{ m}$$

\textbf{Respuestas incorrectas:}
\begin{itemize}
    \item \textbf{A) 25 m:} Divide $50/2$ sin aplicar la fórmula correcta.
    \item \textbf{B) 50 m:} Usa solo la velocidad inicial sin elevarla al cuadrado.
    \item \textbf{D) 250 m:} Olvida el factor 2 en el denominador.
\end{itemize}
\end{respuestabox}

\newpage

\begin{respuestabox}
\subsection{Pregunta 7}
\textcolor{verdecorrecto}{\textbf{Respuesta correcta: B) -5 m/s²}}

\textbf{Justificación:}\\
\textbf{Fórmula utilizada:} Ecuación de velocidad-distancia en MUA (Lección: Movimiento uniformemente acelerado)
$$v_f^2 = v_i^2 + 2ad$$

\textbf{Condición:} Al detenerse, $v_f = 0$

\textbf{Despeje para $a$:}
$$v_f^2 = v_i^2 + 2ad$$
$$0 = v_i^2 + 2ad$$
$$-v_i^2 = 2ad$$
$$a = -\frac{v_i^2}{2d}$$

\textbf{Sustitución:}
$$a = -\frac{(20)^2}{2(40)} = -\frac{400}{80} = -5 \text{ m/s}^2$$

\textbf{Respuestas incorrectas:}
\begin{itemize}
    \item \textbf{A) -10 m/s²:} Divide entre 40 en lugar de 80.
    \item \textbf{C) -2.5 m/s²:} Error en el numerador.
    \item \textbf{D) -0.5 m/s²:} Invierte la fracción.
\end{itemize}
\end{respuestabox}

\begin{respuestabox}
\subsection{Pregunta 8}
\textcolor{verdecorrecto}{\textbf{Respuesta correcta: B) 10 m/s}}

\textbf{Justificación:}\\
\textbf{Fórmula utilizada:} Distancia con velocidad promedio (Lección: Movimiento uniformemente acelerado)
$$d = \frac{v_i + v_f}{2} \times t$$

\textbf{Condición:} $v_f = 2v_0$

\textbf{Despeje para $v_0$:}
$$d = \frac{v_0 + 2v_0}{2} \times t$$
$$d = \frac{3v_0}{2} \times t$$
$$2d = 3v_0 t$$
$$v_0 = \frac{2d}{3t}$$

\textbf{Sustitución:}
$$v_0 = \frac{2(60)}{3(4)} = \frac{120}{12} = 10 \text{ m/s}$$

\textbf{Respuestas incorrectas:}
\begin{itemize}
    \item \textbf{A) 5 m/s:} Error en el cálculo del factor 3.
    \item \textbf{C) 15 m/s:} Divide $60/4$ sin considerar la velocidad promedio.
    \item \textbf{D) 20 m/s:} No considera el factor $\frac{3}{2}$.
\end{itemize}
\end{respuestabox}

\begin{respuestabox}
\subsection{Pregunta 9}
\textcolor{verdecorrecto}{\textbf{Respuesta correcta: C) 80 m}}

\textbf{Justificación:}\\
\textbf{Fórmula utilizada:} Distancia en caída libre (Lección: Movimiento uniformemente acelerado - caída libre)
$$h = \frac{1}{2}gt^2$$

\textbf{Despeje:} Ya está despejada para $h$

\textbf{Datos:} $g = 10$ m/s², $t = 4$ s

\textbf{Sustitución:}
$$h = \frac{1}{2}(10)(4)^2 = 5(16) = 80 \text{ m}$$

\textbf{Respuestas incorrectas:}
\begin{itemize}
    \item \textbf{A) 40 m:} Usa $gt$ en lugar de $\frac{1}{2}gt^2$.
    \item \textbf{B) 60 m:} Error en el cálculo.
    \item \textbf{D) 160 m:} Olvida el factor $\frac{1}{2}$: $(10)(16)$.
\end{itemize}
\end{respuestabox}

\begin{respuestabox}
\subsection{Pregunta 10}
\textcolor{verdecorrecto}{\textbf{Respuesta correcta: B) -10 m/s}}

\textbf{Justificación:}\\
\textbf{Fórmula utilizada:} Velocidad en MRU (Lección: Movimiento rectilíneo uniforme)
$$v = \frac{\Delta x}{\Delta t} = \frac{x_f - x_i}{t_f - t_i}$$

\textbf{Despeje:} Ya está despejada para $v$

\textbf{Sustitución:}
$$v = \frac{10 - 40}{5 - 2} = \frac{-30}{3} = -10 \text{ m/s}$$

El signo negativo indica movimiento en dirección negativa.

\textbf{Respuestas incorrectas:}
\begin{itemize}
    \item \textbf{A) -15 m/s:} Divide entre 2 en lugar de 3.
    \item \textbf{C) -5 m/s:} Error en el cálculo del desplazamiento.
    \item \textbf{D) 10 m/s:} Olvida el signo negativo.
\end{itemize}
\end{respuestabox}

\newpage

\begin{respuestabox}
\subsection{Pregunta 11}
\textcolor{verdecorrecto}{\textbf{Respuesta correcta: C) 150 m}}

\textbf{Justificación:}\\
\textbf{Fórmula utilizada:} Distancia en MRU para ambos móviles (Lección: Movimiento rectilíneo uniforme)
$$d = vt$$

\textbf{Condición de encuentro:} $d_1 = d_2$

Si el segundo móvil tarda tiempo $t$, el primero tarda $t + 4$:

\textbf{Despeje para $t$:}
$$v_1(t + 4) = v_2 t$$
$$v_1 t + 4v_1 = v_2 t$$
$$4v_1 = v_2 t - v_1 t$$
$$4v_1 = t(v_2 - v_1)$$
$$t = \frac{4v_1}{v_2 - v_1}$$

\textbf{Sustitución:}
$$t = \frac{4(15)}{25 - 15} = \frac{60}{10} = 6 \text{ s}$$

\textbf{Cálculo de distancia:}
$$d = v_2 t = 25(6) = 150 \text{ m}$$

\textbf{Respuestas incorrectas:}
\begin{itemize}
    \item \textbf{A) 60 m:} Usa solo $4v_1$.
    \item \textbf{B) 100 m:} Calcula para $t = 4$ s.
    \item \textbf{D) 200 m:} Error en el tiempo o velocidad.
\end{itemize}
\end{respuestabox}
\
\begin{respuestabox}
\subsection{Pregunta 12}
\textcolor{verdecorrecto}{\textbf{Respuesta correcta: C) 32 m}}

\textbf{Justificación:}\\
\textbf{Fórmula utilizada:} Distancia en MUA partiendo del reposo (Lección: Movimiento uniformemente acelerado)
$d = \frac{1}{2}at^2$

\textbf{Paso 1: Encontrar la aceleración}

\textbf{Despeje para $a$:}
$d = \frac{1}{2}at^2$
$2d = at^2$
$a = \frac{2d}{t^2}$

\textbf{Sustitución:}
$a = \frac{2(8)}{(2)^2} = \frac{16}{4} = 4 \text{ m/s}^2$

\textbf{Paso 2: Calcular distancia para $t = 4$ s}

\textbf{Sustitución:}
$d = \frac{1}{2}(4)(4)^2 = 2(16) = 32 \text{ m}$

\textbf{Respuestas incorrectas:}
\begin{itemize}
    \item \textbf{A) 16 m:} Duplica la distancia sin considerar $d \propto t^2$.
    \item \textbf{B) 24 m:} Triplica incorrectamente.
    \item \textbf{D) 64 m:} Usa $t^3$ en lugar de $t^2$.
\end{itemize}
\end{respuestabox}
\newpage

\begin{respuestabox}
\subsection{Pregunta 13}
\textcolor{verdecorrecto}{\textbf{Respuesta correcta: C) 90 m}}

\textbf{Justificación:}\\
\textbf{Paso 1: Calcular tiempo de caída} (Lección: Movimiento uniformemente acelerado - caída libre)
$h = \frac{1}{2}gt^2$

\textbf{Despeje para $t$:}
$2h = gt^2$
$t^2 = \frac{2h}{g}$
$t = \sqrt{\frac{2h}{g}}$

\textbf{Sustitución:}
$t = \sqrt{\frac{2(45)}{10}} = \sqrt{\frac{90}{10}} = \sqrt{9} = 3 \text{ s}$

\textbf{Paso 2: Calcular distancia horizontal} (Lección: Movimiento rectilíneo uniforme)
$x = v_x t$

\textbf{Sustitución:}
$x = 30(3) = 90 \text{ m}$

\textbf{Respuestas incorrectas:}
\begin{itemize}
    \item \textbf{A) 30 m:} Usa $t = 1$ s.
    \item \textbf{B) 60 m:} Usa $t = 2$ s.
    \item \textbf{D) 135 m:} Error en el cálculo del tiempo.
\end{itemize}
\end{respuestabox}


\begin{respuestabox}
\subsection{Pregunta 14}
\textcolor{verdecorrecto}{\textbf{Respuesta correcta: B) 1.5 m/s²}}

\textbf{Justificación:}\\
\textbf{Concepto:} En una gráfica v-t, la pendiente es la aceleración (Lección: Movimiento uniformemente acelerado - interpretación gráfica)

\textbf{Fórmula:}
$a = \frac{\Delta v}{\Delta t}$

\textbf{Despeje:} Ya está despejada para $a$

\textbf{Datos de la gráfica:} Parte de reposo ($v_i = 0$), llega a $v_f = 12$ m/s en $t = 8$ s

\textbf{Sustitución:}
$a = \frac{12 - 0}{8} = \frac{12}{8} = 1.5 \text{ m/s}^2$

\textbf{Respuestas incorrectas:}
\begin{itemize}
    \item \textbf{A) 0.67 m/s²:} Invierte la división.
    \item \textbf{C) 3 m/s²:} Divide entre 4 en lugar de 8.
    \item \textbf{D) 6 m/s²:} Divide entre 2.
\end{itemize}
\end{respuestabox}

\begin{respuestabox}
\subsection{Pregunta 15}
\textcolor{verdecorrecto}{\textbf{Respuesta correcta: C) 30 m/s}}

\textbf{Justificación:}\\
\textbf{Concepto:} El tiempo de subida es igual al tiempo de bajada (Lección: Movimiento uniformemente acelerado - caída libre)

\textbf{Tiempo de subida:}
$t_{subida} = \frac{t_{total}}{2} = \frac{6}{2} = 3 \text{ s}$

\textbf{Fórmula:} En la altura máxima, $v_f = 0$
$v_f = v_i - gt$

\textbf{Despeje para $v_i$:}
$0 = v_i - gt$
$v_i = gt$

\textbf{Sustitución:}
$v_i = 10(3) = 30 \text{ m/s}$

\textbf{Respuestas incorrectas:}
\begin{itemize}
    \item \textbf{A) 15 m/s:} Usa $t = 1.5$ s.
    \item \textbf{B) 20 m/s:} Usa $t = 2$ s.
    \item \textbf{D) 60 m/s:} Usa el tiempo total en lugar del tiempo de subida.
\end{itemize}
\end{respuestabox}
